% GENERATED FILE. Edit canonical lessons, then run npm run generate:books.
% canonical-source-hash: fnv1a64:463e9d365904772b
% canonical-lessons: LA-C327-sponsa, LA-C327-sponsalia, LA-C327-despondeo, LA-C327-nupta, LA-C327-matrimonium

\chapter{Spōnsālia}
\label{ch:la-sponsalia}
\hlchaptermodality{327}

\begin{chapteropening}
\textbf{By the end of this chapter:} \emph{I can talk about getting engaged and married: the bride-to-be, the betrothal, to promise in marriage, the bride and marriage.}

Complete the last lesson of chapter 327: I can talk about getting engaged and married: the bride-to-be, the betrothal, to promise in marriage, the bride and marriage.
\end{chapteropening}

\section[spōnsa, spōnsae]{spōnsa, spōnsae — a bride-to-be, a fiancée}

\label{lesson:LA-C327-sponsa}



\begin{quote}
Before the new one: say the Latin for orphaned, then the Latin for a ward.
\end{quote}

\subsection*{You'll want to know: spōnsa, spōnsae}
\textbf{spōnsa, spōnsae} — \textquotedblleft{}a bride-to-be, a fiancée\textquotedblright{}.

\begin{quote}
Spōnsa eius ex oppidō vīcīnō erat.
\end{quote}

\begin{quote}
\emph{His bride-to-be was from a neighbouring town.}
\end{quote}

Feminine, a participle: a girl promised. Her man is the \emph{spōnsus}. English \emph{spouse}.

\begin{cousinweb}
\textbf{spōnsa} is the feminine participle of \emph{spondeō}, \textquotedblleft{}to promise solemnly\textquotedblright{} (EtymDB): she who has been promised. English \emph{spouse} comes from it through French.
\end{cousinweb}

\begin{grammarlens}[title={What you've built}]
One more word for telling, explaining and describing at B1.
\end{grammarlens}

\subsection*{Your turn}
\begin{itemize}
\raggedright
  \item \emph{Say it:} \emph{spōnsa, spōnsae}
  \item \emph{Say it:} \emph{spōnsa, spōnsae}, once more
  \item \emph{Say it:} a sentence of your own with \emph{spōnsa, spōnsae}
  \item \emph{Recall:} say the Latin for orphaned, then the Latin for a ward, then say \emph{spōnsa, spōnsae} again
\end{itemize}

\begin{tcolorbox}[breakable,colback=teal!4,colframe=teal!35!black,title={Before you move on}]
What does spōnsa, spōnsae mean? (\textquotedblleft{}A bride-to-be, a fiancée\textquotedblright{}.) Say it once more.
\end{tcolorbox}
\section[spōnsālia, spōnsālium]{spōnsālia, spōnsālium — a betrothal; an engagement party}

\label{lesson:LA-C327-sponsalia}



\begin{quote}
Before the new one: say the Latin for a ward, then the Latin for a bride-to-be.
\end{quote}

\subsection*{You'll want to know: spōnsālia, spōnsālium}
\textbf{spōnsālia, spōnsālium} — \textquotedblleft{}a betrothal; an engagement party\textquotedblright{}.

\begin{quote}
Ad spōnsālia multī amīcī vēnērunt.
\end{quote}

\begin{quote}
\emph{Many friends came to the engagement party.}
\end{quote}

Neuter, plural only. A ring was often given at the \emph{spōnsālia}.

\begin{cousinweb}
\textbf{spōnsālia} is built on \emph{spōnsālis}, \textquotedblleft{}of a betrothal\textquotedblright{}, from \emph{spōnsa} (EtymDB).
\end{cousinweb}

\begin{grammarlens}[title={What you've built}]
One more word for telling, explaining and describing at B1.
\end{grammarlens}

\subsection*{Your turn}
\begin{itemize}
\raggedright
  \item \emph{Say it:} \emph{spōnsālia, spōnsālium}
  \item \emph{Say it:} \emph{spōnsālia, spōnsālium}, once more
  \item \emph{Say it:} a sentence of your own with \emph{spōnsālia, spōnsālium}
  \item \emph{Recall:} say the Latin for a ward, then the Latin for a bride-to-be, then say \emph{spōnsālia, spōnsālium} again
\end{itemize}

\begin{tcolorbox}[breakable,colback=teal!4,colframe=teal!35!black,title={Before you move on}]
What does spōnsālia, spōnsālium mean? (\textquotedblleft{}A betrothal; an engagement party\textquotedblright{}.) Say it once more.
\end{tcolorbox}
\section[dēspondeō, dēspondēre]{dēspondeō, dēspondēre — to promise in marriage, to betroth}

\label{lesson:LA-C327-despondeo}



\begin{quote}
Before the new one: say the Latin for a bride-to-be, then the Latin for a betrothal.
\end{quote}

\subsection*{You'll want to know: dēspondeō, dēspondēre}
\textbf{dēspondeō, dēspondēre} — \textquotedblleft{}to promise in marriage, to betroth\textquotedblright{}.

\begin{quote}
Pater fīliam suam amīcō dēspondit.
\end{quote}

\begin{quote}
\emph{The father promised his daughter in marriage to a friend.}
\end{quote}

\emph{dēspondeō, dēspondēre, dēspondī, dēspōnsum}. \emph{Animum dēspondēre}: to lose heart, which gives English \emph{despond}.

\begin{cousinweb}
\textbf{dēspondeō} is built on \emph{spondeō}, \textquotedblleft{}to promise solemnly\textquotedblright{}, with \emph{dē-} (EtymDB): to promise away. English \emph{despond} comes from its use with \emph{animum}, \textquotedblleft{}to lose heart\textquotedblright{}.
\end{cousinweb}

\begin{grammarlens}[title={What you've built}]
One more word for telling, explaining and describing at B1.
\end{grammarlens}

\subsection*{Your turn}
\begin{itemize}
\raggedright
  \item \emph{Say it:} \emph{dēspondeō, dēspondēre}
  \item \emph{Say it:} \emph{dēspondeō, dēspondēre}, once more
  \item \emph{Say it:} a sentence of your own with \emph{dēspondeō, dēspondēre}
  \item \emph{Recall:} say the Latin for a bride-to-be, then the Latin for a betrothal, then say \emph{dēspondeō, dēspondēre} again
\end{itemize}

\begin{tcolorbox}[breakable,colback=teal!4,colframe=teal!35!black,title={Before you move on}]
What does dēspondeō, dēspondēre mean? (\textquotedblleft{}To promise in marriage, to betroth\textquotedblright{}.) Say it once more.
\end{tcolorbox}
\section[nūpta, nūptae]{nūpta, nūptae — a bride, a married woman}

\label{lesson:LA-C327-nupta}



\begin{quote}
Before the new one: say the Latin for a betrothal, then the Latin for to promise in marriage.
\end{quote}

\subsection*{You'll want to know: nūpta, nūptae}
\textbf{nūpta, nūptae} — \textquotedblleft{}a bride, a married woman\textquotedblright{}.

\begin{quote}
Nūpta ad domum marītī ducta est.
\end{quote}

\begin{quote}
\emph{The bride was led to her husband's house.}
\end{quote}

Feminine, the participle of \emph{nūbō} (taught), said of the woman who marries. \emph{Nova nūpta}: the new bride.

\begin{cousinweb}
\textbf{nūpta} is the feminine participle of \emph{nūbō}, \textquotedblleft{}to marry\textquotedblright{} (EtymDB), which is inherited from a Proto-Indo-European root meaning \textquotedblleft{}to wed\textquotedblright{}. English \emph{nuptial} comes from the same family.
\end{cousinweb}

\begin{grammarlens}[title={What you've built}]
One more word for telling, explaining and describing at B1.
\end{grammarlens}

\subsection*{Your turn}
\begin{itemize}
\raggedright
  \item \emph{Say it:} \emph{nūpta, nūptae}
  \item \emph{Say it:} \emph{nūpta, nūptae}, once more
  \item \emph{Say it:} a sentence of your own with \emph{nūpta, nūptae}
  \item \emph{Recall:} say the Latin for a betrothal, then the Latin for to promise in marriage, then say \emph{nūpta, nūptae} again
\end{itemize}

\begin{tcolorbox}[breakable,colback=teal!4,colframe=teal!35!black,title={Before you move on}]
What does nūpta, nūptae mean? (\textquotedblleft{}A bride, a married woman\textquotedblright{}.) Say it once more.
\end{tcolorbox}
\section[mātrimōnium, mātrimōniī]{mātrimōnium, mātrimōniī — marriage}

\label{lesson:LA-C327-matrimonium}



\begin{quote}
Before the new one: say the Latin for to promise in marriage, then the Latin for a bride.
\end{quote}

\subsection*{You'll want to know: mātrimōnium, mātrimōniī}
\textbf{mātrimōnium, mātrimōniī} — \textquotedblleft{}marriage\textquotedblright{}.

\begin{quote}
Puellam in mātrimōnium dūxit.
\end{quote}

\begin{quote}
\emph{He married the girl.}
\end{quote}

Neuter. \emph{In mātrimōnium dūcere}: (of a man) to marry. English \emph{matrimony}.

\begin{cousinweb}
\textbf{mātrimōnium} is built on \emph{māter}, \textquotedblleft{}mother\textquotedblright{}, with the ending \emph{-mōnium} (EtymDB): the state of a mother, of a wife. English \emph{matrimony} comes from it.
\end{cousinweb}

\begin{grammarlens}[title={What you've built}]
One more word for telling, explaining and describing at B1.
\end{grammarlens}

\subsection*{Your turn}
\begin{itemize}
\raggedright
  \item \emph{Say it:} \emph{mātrimōnium, mātrimōniī}
  \item \emph{Say it:} \emph{mātrimōnium, mātrimōniī}, once more
  \item \emph{Say it:} a sentence of your own with \emph{mātrimōnium, mātrimōniī}
  \item \emph{Recall:} say the Latin for to promise in marriage, then the Latin for a bride, then say \emph{mātrimōnium, mātrimōniī} again
\end{itemize}

\begin{tcolorbox}[breakable,colback=teal!4,colframe=teal!35!black,title={Before you move on}]
What does mātrimōnium, mātrimōniī mean? (\textquotedblleft{}Marriage\textquotedblright{}.) Say it once more.
\end{tcolorbox}
